% ---
% id: foundations/lem:floor-stricter
% kind: lem
% title: "The positivity floor against entrywise closeness to the cone"
% summary: >-
%   A Hermitian matrix with least eigenvalue at least −δ lies within entrywise
%   distance δ of the positive cone, but entrywise distance δ only guarantees a
%   least eigenvalue of −dδ in dimension d, and that factor is attained. The two
%   tests are different quantities and never interchangeable.
% prominence: standard
% status: accepted
% uses: []
% mentions: [def:tolerance]
% assumes: []
% refs: []
% provenance: newly-derived
% lean:
% ---

The declared tolerance governs two kinds of test, entrywise defects of
differences and a floor on the least eigenvalue
(Definition~\ref{def:tolerance}). The lemma compares them.

\begin{lemma}\label{lem:floor-stricter}
Let $d\ge1$ and $\delta\ge0$, and write $\opn{\cdot}$ for the operator norm and
$\mxn{\cdot}$ for the entrywise maximum modulus of complex $d\times d$
matrices.
\begin{enumerate}[label=(\roman*),leftmargin=2.4em]
\item $\mxn{Y}\le\opn{Y}\le d\,\mxn{Y}$ for every complex $d\times d$
  matrix $Y$.
\item If $X$ is Hermitian and $\lmin(X)\ge-\delta$, then
  $X+\delta\,\idon{d}\succeq0$ and
  $\mxn{(X+\delta\,\idon{d})-X}=\delta$.
\item If $X$ is Hermitian and $\mxn{X-P}\le\delta$ for some $P\succeq0$, then
  $\lmin(X)\ge-d\delta$. The bound is attained: for the $d\times d$ matrix $E$
  with every entry equal to $1$, the matrices $X\coloneqq-\delta E$ and
  $P\coloneqq0$ satisfy $\mxn{X-P}=\delta$ and $\lmin(X)=-d\delta$.
\end{enumerate}
\end{lemma}

\begin{proof}
(i) For standard basis vectors $e_i,e_l$,
$|Y_{il}|=|\langle e_i,Ye_l\rangle|\le\opn{Y}$ by the Cauchy--Schwarz
inequality, which is the first inequality. For a unit vector $v$, the
Cauchy--Schwarz inequality applied to each row gives
\begin{equation*}
  \|Yv\|^{2}=\sum_i\Big|\sum_lY_{il}v_l\Big|^{2}
  \le\sum_i\Big(\sum_l|Y_{il}|^{2}\Big)\Big(\sum_l|v_l|^{2}\Big)
  =\sum_{i,l}|Y_{il}|^{2}\le d^{2}\mxn{Y}^{2},
\end{equation*}
so $\opn{Y}\le d\,\mxn{Y}$.

(ii) The Hermitian matrix $X+\delta\,\idon{d}$ has the eigenvalues of $X$
shifted by $\delta$, which are non-negative, so it is positive semidefinite.
The difference $\delta\,\idon{d}$ has diagonal entries $\delta$ and
off-diagonal entries $0$, so its entrywise maximum modulus is $\delta$.

(iii) $X-P$ is Hermitian, $X$ and $P$ being so. For a unit vector $v$,
\begin{equation*}
  \langle v,Xv\rangle=\langle v,Pv\rangle+\langle v,(X-P)v\rangle
  \ge0-\opn{X-P}\ge-d\,\mxn{X-P}\ge-d\delta,
\end{equation*}
by $P\succeq0$, the Cauchy--Schwarz inequality and (i). Since
$\lmin(X)=\min_{\|v\|=1}\langle v,Xv\rangle$ for Hermitian $X$, this gives
$\lmin(X)\ge-d\delta$. For the example, let $u$ be the unit vector with every
entry $d^{-1/2}$; then $E=d\,uu^{\dagger}$, whose eigenvalues are $d$, on $u$,
and $0$, on the orthogonal complement of $u$. So the least eigenvalue of
$-\delta E$ is $-d\delta$, and every entry of $-\delta E-0$ has modulus
$\delta$.
\end{proof}
